% Post-CP32 candidate. Complete direct proof; independent written audit pending.
\subsection{The exact group structure left by full hyperplane inertia}
\label{subsec:full-radical-image-structure}

Throughout this subsection \(V\) is an \(r\)-dimensional vector space
over \(\mathbf F_2\), with \(r\geq2\). For a nonzero
\(\ell\in V^\vee\), put
\[
 U_\ell=\{1+v\ell:v\in\ker\ell\}\subset\operatorname{GL}(V).
 \tag{FR1}
\]
Here \(v\ell\) denotes the endomorphism \(x\mapsto\ell(x)v\).
The condition \(\ell(v)=0\) makes every displayed map invertible,
and \(U_\ell\) is the full group of order \(2^{r-1}\).
All stabilizers described as pointwise below act as the identity on
every vector of the specified subspace.

\begin{lemma}[Full radicals generate a complete pointwise stabilizer]
\label{lem:root-full-radicals-pointwise-stabilizer}
\status{Candidate; complete direct proof supplied for independent audit.}
Let \(S\subset V^\vee\setminus\{0\}\) be nonempty. Set
\[
 W=\langle S\rangle_{\mathbf F_2},\qquad K=W^\perp,
 \qquad m=\dim W=\dim(V/K).
 \tag{FR2}
\]
Then
\[
 \langle U_\ell:\ell\in S\rangle
   =\{g\in\operatorname{GL}(V):g|_K=1_K\}.
 \tag{FR3}
\]
Its fixed subspace in \(V\) is exactly \(K\).
\end{lemma}
\begin{proof}
Every generator fixes \(K\) pointwise, giving one inclusion.
Choose a basis \(\ell_1,\ldots,\ell_m\) of \(W\) from \(S\).
Choose a complement \(X\) to \(K\) and a basis
\(x_1,\ldots,x_m\) of \(X\) dual to these functionals.
Thus \(\ell_i(x_j)=\delta_{ij}\) and \(\ell_i(K)=0\).

For every \(k\in K\), \(U_{\ell_i}\) contains
\(1+k\ell_i\). These elements generate all transformations
\[
 \begin{pmatrix}1_K&B\\0&1_X\end{pmatrix},
 \qquad B\in\operatorname{Hom}_{\mathbf F_2}(X,K),
 \tag{FR4}
\]
in the ordered decomposition \(V=K\oplus X\).
For \(i\ne j\), it also contains \(1+x_j\ell_i\).
These act trivially on \(K\) and are the elementary row additions
on \(X\). If \(m\geq2\), they generate
\(\operatorname{GL}(X)\) by Gaussian elimination; a row interchange
is three row additions and over \(\mathbf F_2\) no nontrivial
diagonal scaling is needed. If \(m=1\), this general linear group
is trivial, so the same assertion holds without any row additions.

Every automorphism fixing \(K\) pointwise has the block form
\(\left(\begin{smallmatrix}1&B\\0&A\end{smallmatrix}\right)\)
with \(A\in\operatorname{GL}(X)\), and is a product of the
two kinds just constructed. This proves (FR3).
If \(K\ne0\), any vector outside \(K\) is moved by some
shear (FR4). If \(K=0\), then \(m=r\geq2\), and
\(\operatorname{GL}(V)\) fixes no nonzero vector. Thus the fixed
space is \(K\) in both cases. The hypothesis \(r\geq2\)
excludes the exceptional trivial action on a one-dimensional space.
\end{proof}

\begin{proposition}[A group containing one full radical]
\label{prop:root-full-radical-parabolic-form}
\status{Candidate; complete direct proof supplied for independent audit.}
Suppose \(G\subseteq\operatorname{GL}(V)\) contains \(U_{\ell_0}\).
Let \(N\) be its normal closure in \(G\), and put \(K=V^N\).
Then \(K\) is a proper \(G\)-stable subspace, and \(N\) is
its complete pointwise stabilizer in \(\operatorname{GL}(V)\).
If \(H\) is the image of the restriction homomorphism
\(G\to\operatorname{GL}(K)\), then, for any complement \(X\),
\[
 G=\left\{
   \begin{pmatrix}h&B\\0&A\end{pmatrix}:
   h\in H,\ A\in\operatorname{GL}(X),\
   B\in\operatorname{Hom}_{\mathbf F_2}(X,K)
       \right\}.
 \tag{FR5}
\]
Equivalently this is the full inverse image of \(H\) under the
restriction map from the setwise stabilizer of \(K\) to
\(\operatorname{GL}(K)\). In particular,
\[
 G=\operatorname{Hom}(V/K,K)\rtimes
        \bigl(H\times\operatorname{GL}(V/K)\bigr),
 \tag{FR6}
\]
after the chosen splitting; the displayed equality identifies the
natural block action, not merely an abstract group isomorphism.
\end{proposition}
\begin{proof}
Conjugation gives
\(gU_\ell g^{-1}=U_{\ell\circ g^{-1}}\).
Apply the lemma to the orbit
\(S=\{\ell_0\circ g^{-1}:g\in G\}\).
It shows that \(N\) is the complete pointwise stabilizer of
\(K_0=\bigcap_{\ell\in S}\ker\ell\), and that \(V^N=K_0\).
Thus \(K=K_0\) is proper and is \(G\)-stable, since \(N\)
is normal. The restriction homomorphism from its setwise
stabilizer to \(\operatorname{GL}(K)\) is surjective and has
kernel \(N\). A subgroup containing this entire kernel is
the full inverse image of its image. This proves (FR5).
The block-diagonal lifts of \(H\times\operatorname{GL}(X)\)
and the shears give (FR6); their action on a shear is
\(B\mapsto hBA^{-1}\).
\end{proof}

\begin{corollary}[Submodules and the one large nonzero orbit]
\label{cor:root-full-radical-submodules-orbits}
\status{Candidate; complete direct proof supplied for independent audit.}
In the proposition, \(K\) is the unique maximal proper
\(G\)-submodule of \(V\). The group acts transitively on
\(V\setminus K\), an orbit of cardinality
\(2^r-2^{\dim K}\); its orbits inside \(K\) are precisely
the \(H\)-orbits. If \((V^\vee)^G=0\), then
\(\dim(V/K)\geq2\). Conversely, if \(\dim(V/K)\geq2\),
then \((V^\vee)^G=0\), even when \(K\ne0\).
The following conditions are equivalent:
\[
 K=0,\qquad V\text{ is irreducible as a }G\text{-module},
 \qquad G\text{ is transitive on }V\setminus\{0\},
 \qquad G=\operatorname{GL}(V).
 \tag{FR7}
\]
\end{corollary}
\begin{proof}
The full general linear group on \(X\) sends any nonzero
quotient coordinate to any other, and a shear then assigns
the \(K\)-coordinate arbitrarily. This proves transitivity
outside \(K\). Formula (FR5) proves the assertion inside it.
If a \(G\)-stable subspace contains a vector outside \(K\),
its shear differences contain every vector of \(K\), and
its image in \(V/K\) is the whole quotient. It is therefore
\(V\). Every proper submodule is consequently contained in
\(K\), proving its uniqueness, including \(K=0\).

If \(m=1\), the quotient \(V\to V/K\simeq\mathbf F_2\)
is a nonzero invariant covector. Suppose \(m\geq2\).
When \(K\ne0\), invariance under every shear forces any
invariant covector to vanish on \(K\); when \(K=0\) this
condition is automatic. The induced covector on \(X\) is
fixed by its full general linear group and hence is zero.
This proves both covector assertions. Finally (FR5) gives
the full group when \(K=0\), and a nonzero proper \(K\)
prevents both irreducibility and transitivity on nonzero
vectors. This proves (FR7).
\end{proof}

\begin{corollary}[What remains for every fixed odd principal-cube system]
\label{cor:root-odd-image-parabolic-alternative}
\status{Candidate; uses the separately proved full-inertia theorem.}
Let \((A,\Theta)\) be a fixed nonordinary principally polarized
abelian surface over a perfect field of characteristic two, and
let \(d\geq1\) be odd. For the geometric relative Lang image
\(G_d\) on its rank-\(r=9d^2-1\) space, there are a canonical
proper submodule \(K_d\) and a subgroup
\(H_d\subseteq\operatorname{GL}(K_d)\) such that its natural
action is exactly (FR5), with \(\dim(V/K_d)\geq2\).
After base change to an algebraic closure, its nonzero Lang cover
has one geometric connected component corresponding to
\(V\setminus K_d\), of degree \(2^r-2^{\dim K_d}\), and
the remaining geometric connected components have the degrees of the nonzero
\(H_d\)-orbits in \(K_d\).
\end{corollary}
\begin{proof}
Corollary~\ref{cor:fixed-ppav-odd-full-hyperplane-inertia}
supplies an actual full \(U_\ell\) in \(G_d\).
Theorem~\ref{thm:nonordinary-ppav-all-covectors} gives
\((V^\vee)^{G_d}=0\) for these odd systems. Apply the
proposition and preceding corollary. The resulting \(K_d\)
is canonical because it is the unique maximal proper
submodule; it does not depend on the initially chosen full
radical. After the same geometric base change, connected components of a
finite etale cover over a connected normal base correspond to orbits of its
fundamental group on the geometric fibre. Applying this
elementary equivalence to the nonzero fixed-vector cover
gives the stated degrees.
\end{proof}

\par\smallskip\noindent\textbf{Scope.}
The proposition identifies the entire form of any remaining
proper image once \(K_d\) and its action \(H_d\) are known.
It does not determine \(K_d\), and therefore does not supply
a new exact full-image degree or settle general A.4.
The absence of invariant covectors alone cannot exclude the
proper alternative, as its explicit covector computation shows.
These are direct finite-dimensional proofs; no numerical
experiment or external classification is used. The normal
closure of the known full radicals is not silently replaced
by the group generated by all divisorial inertia.
